\documentclass[12pt]{article}

% --- Packages ---
\usepackage[margin=1in]{geometry}
\usepackage{amsmath,amssymb,amsthm}
\usepackage{mathtools}
\usepackage{graphicx}
\usepackage{booktabs}
\usepackage{multirow}
\usepackage{natbib}
\usepackage{hyperref}
\usepackage{cleveref}
\usepackage{algorithm}
\usepackage{algpseudocode}
\usepackage{bm}

% --- Theorem environments ---
\newtheorem{theorem}{Theorem}
\newtheorem{lemma}[theorem]{Lemma}
\newtheorem{proposition}[theorem]{Proposition}
\newtheorem{corollary}[theorem]{Corollary}
\theoremstyle{definition}
\newtheorem{definition}[theorem]{Definition}
\newtheorem{example}[theorem]{Example}
\theoremstyle{remark}
\newtheorem{remark}[theorem]{Remark}

% --- Notation shortcuts ---
\newcommand{\E}{\mathbb{E}}
\newcommand{\Var}{\mathrm{Var}}
\newcommand{\Rset}{\mathbb{R}}
\newcommand{\bone}{\mathbf{1}}
\newcommand{\balpha}{\boldsymbol{\alpha}}
\newcommand{\bbeta}{\boldsymbol{\beta}}
\newcommand{\blam}{\boldsymbol{\lambda}}
\newcommand{\btheta}{\bm{\theta}}
\DeclareMathOperator*{\argmax}{arg\,max}
\DeclareMathOperator*{\argmin}{arg\,min}

% --- Title ---
\title{Weibull Component Estimation from $k$-out-of-$n$ System Data:\\
  Why Exponential Results Don't Generalize}
\author{Alexander Towell\thanks{Department of Computer Science,
  Southern Illinois University Edwardsville.
  Email: \texttt{lex@metafunctor.com}.
  ORCID: 0000-0001-6443-9897.}}
\date{}

\begin{document}
\maketitle

\begin{abstract}
Exponential component lifetime estimation from system data is a
well-understood problem: closed-form likelihoods via
inclusion-exclusion, analytical EM algorithms, and provable
identifiability results are available for parallel systems with
moderate numbers of components.  A natural next step is to extend
these results to the Weibull family, which captures the increasing
and decreasing failure rates that arise in wear-out and early-life
failure.  We show that this extension is \emph{not} a minor
generalization but a fundamentally harder problem.  The additional
shape parameters double the parameter count, destroy the
inclusion-exclusion closed forms that underpin exponential
tractability, degrade optimizer convergence from near-perfect to as
low as 47\%, inflate root mean squared error by factors of 10--30,
and---most strikingly---change the optimal system structure for
estimation in a shape-dependent way.  Under constant hazard
(exponential baseline), the optimal $k$ for a $k$-out-of-5 system
is $k^* = 4$; under increasing hazard (shape $= 2$), it shifts to
$k^* = 1$; under mixed shapes, $k^* = 2$.  Series systems
($k = n$) are catastrophic for Weibull estimation, with RMSE
exceeding 360 under mixed shapes.  We develop and implement an EM
algorithm whose E-step computes truncated Weibull moments via
incomplete gamma functions and whose M-step uses profile
optimization over component shapes; the EM achieves substantially
higher convergence rates than direct MLE, particularly for DFR
components.  We derive the periodic inspection (Scheme~1) likelihood
and demonstrate empirically that it resolves the shape--scale
ambiguity inherent in system-level observation: shape RMSE drops by
an order of magnitude, with the largest gains for deeply-censored
fast components.  Monte Carlo experiments with two- to
five-component systems across shape scenarios spanning DFR through
IFR quantify the severity of the exponential-to-Weibull transition
and validate the complete estimation framework.
\end{abstract}

\noindent\textbf{Keywords:} Weibull distribution,
$k$-out-of-$n$ systems, system reliability, maximum likelihood,
EM algorithm, periodic inspection, component estimation

% ===================================================================
\section{Introduction}
\label{sec:introduction}
% ===================================================================

The exponential distribution occupies a privileged position in
system reliability analysis.  Its memoryless property implies that
the minimum of independent exponentials is itself exponential, the
hazard rate of a series system is the sum of component hazard
rates, and inclusion-exclusion expansions of the parallel system CDF
yield closed-form integrals amenable to maximum likelihood
estimation \citep{Towell2026}.  These analytical conveniences have
enabled a thorough understanding of component estimation from
system-level data in the exponential case: identifiability results
\citep{BasuGhosh1980}, closed-form EM algorithms, and sharp
characterizations of the information loss due to system-level
observation \citep{Towell2026}.

In practice, however, the constant-hazard assumption is often
untenable.  Mechanical wear, fatigue, and degradation produce
increasing failure rates; manufacturing defects and infant mortality
produce decreasing failure rates; and systems with subpopulations of
strong and weak components exhibit bathtub-shaped hazard curves.
The Weibull distribution \citep{Weibull1951}, with its shape
parameter $\alpha$ governing the hazard trend and its scale
parameter $\beta$ governing the characteristic life, is the standard
parametric model for these phenomena \citep{MeekerEscobar1998,
LawlessBook2003}.

Given the centrality of the Weibull family in reliability
engineering, extending the exponential system estimation results to
Weibull components is a natural and important goal.  The purpose of
this paper is to show that this extension is \emph{fundamentally
harder} than the exponential case---not merely more computationally
expensive, but qualitatively different in its structure,
identifiability landscape, and practical viability.

\subsection*{Why exponential doesn't generalize}

The difficulty is not a single obstacle but a constellation of
interconnected complications:

\begin{enumerate}
\item \textbf{Parameter count.}
  An exponential system with $m$ components has $m$ parameters
  (one rate per component).  A Weibull system has $2m$ parameters
  ($m$ shapes plus $m$ scales).  With $m = 5$, the estimation
  problem jumps from 5 to 10 parameters, and the Fisher information
  per parameter thins correspondingly.

\item \textbf{Loss of closed-form integrals.}
  The exponential system CDF
  $F_T(t) = \prod_j (1 - e^{-\lambda_j t})$ admits an
  inclusion-exclusion expansion into a sum of exponentials, each of
  which integrates in closed form.  The Weibull system CDF
  $F_T(t) = \prod_j (1 - e^{-(t/\beta_j)^{\alpha_j}})$
  has no such expansion.  Each factor
  $1 - e^{-(t/\beta_j)^{\alpha_j}}$ is \emph{not} a linear
  function of exponentials unless $\alpha_j = 1$, so
  the product cannot be decomposed into a finite signed sum of
  elementary terms.  All integrals required for censored
  likelihoods and EM E-steps must be computed numerically.

\item \textbf{Series catastrophe.}
  For exponential components in a series system, the system hazard
  rate $h_T(t) = \sum_j \lambda_j$ is a sum of constants, so the
  system is itself exponential with rate $\Lambda = \sum_j \lambda_j$.
  Individual rates are not identifiable, but $\Lambda$ is trivially
  estimated.  For Weibull components in a series system, the system
  hazard rate $h_T(t) = \sum_j (\alpha_j / \beta_j)(t/\beta_j)^{\alpha_j - 1}$
  is a sum of power functions with different exponents.  Individual
  $(\alpha_j, \beta_j)$ pairs are not identifiable from $T_{\min}$
  alone, and the non-constant hazard creates a multi-modal
  likelihood landscape that numerical optimizers cannot navigate
  reliably.

\item \textbf{Shape-dependent optimal structure.}
  In the exponential case, the optimal system structure ($k$ in a
  $k$-out-of-$n$ system) for component estimation depends on the
  rate heterogeneity but not on a shape parameter---because there
  is none.  In the Weibull case, the shape parameter fundamentally
  alters which system structure is optimal: increasing hazard favors
  parallel systems ($k = 1$), while constant hazard favors
  near-series structures.  A practitioner who designs a test system
  based on exponential intuition will choose the wrong $k$ when
  components have non-unit shapes.
\end{enumerate}

\subsection*{Contributions}

We make the following contributions:

\begin{enumerate}
\item \textbf{Numerical likelihood framework} (\Cref{sec:likelihood}).
  We derive the system density for heterogeneous Weibull components
  in arbitrary coherent systems via critical-state decomposition,
  and show that the likelihood is well-defined and continuously
  differentiable but admits no closed-form simplification beyond
  the exponential special case.

\item \textbf{Quantification of the exponential--Weibull gap}
  (\Cref{sec:simulations}).
  Through systematic Monte Carlo experiments with $m = 5$ Weibull
  components across three shape scenarios (constant, increasing,
  mixed) and five $k$-out-of-5 structures, we quantify the
  degradation: RMSE inflates by factors of 10--30, convergence
  drops to 47--93\%, and series systems become completely unviable
  (RMSE $> 360$).

\item \textbf{Shape-dependent optimal $k$}
  (\Cref{sec:sim-optimal-k}).
  We show that the optimal $k$ for estimation quality depends on
  the shape parameter in a systematic way: $k^* = 4$ for constant
  hazard, $k^* = 1$ for increasing hazard, and $k^* = 2$ for mixed
  shapes.  This is the first evidence that system design for Weibull
  estimation cannot be guided by exponential theory.

\item \textbf{EM algorithm for Weibull systems}
  (\Cref{sec:em-algorithm,sec:em-results}).
  We develop and implement an EM algorithm for Weibull component
  estimation from parallel system data.  The E-step computes
  truncated Weibull moments via incomplete gamma functions; the
  M-step uses profile optimization over each component's shape
  parameter with closed-form scale update.  Empirically, the EM
  achieves near-perfect convergence across all shape
  configurations tested, substantially outperforming the direct MLE
  for DFR and mixed-shape scenarios.

\item \textbf{Periodic inspection likelihood and validation}
  (\Cref{sec:scheme1,sec:scheme1-results}).
  We derive the likelihood for Scheme~1 (periodic inspection) and
  validate it empirically.  Scheme~1 resolves the shape--scale
  ambiguity inherent in system-level observation: shape RMSE drops
  by an order of magnitude relative to Scheme~0, with the largest
  gains for deeply-censored fast components.  Even coarse
  inspection intervals provide substantial improvement.
\end{enumerate}

\subsection*{Relation to companion work}

This paper is the sequel to \citet{Towell2026}, which established
the censoring equivalence for exponential parallel systems, derived
closed-form MLEs via inclusion-exclusion, and proved
identifiability via reversed hazard rate tail extraction.  The
present paper takes the natural next step---extending to the
Weibull family---and shows that this step is far from routine.  The
exponential results of \citet{Towell2026} serve as both the
baseline for comparison and the theoretical foundation upon which
the Weibull framework is built.

\subsection*{Outline}

\Cref{sec:model} introduces the $k$-out-of-$n$ system model with
Weibull components and reviews the censoring interpretation.
\Cref{sec:likelihood} develops the system density and likelihood
via critical-state decomposition, emphasizing the loss of
closed-form structure.
\Cref{sec:why-hard} explains why exponential results do not
generalize, contrasting the mathematical structure of the two
families point by point.
\Cref{sec:simulations} reports the Monte Carlo study comparing
Weibull and exponential estimation across system structures and
shape scenarios via direct MLE.
\Cref{sec:em-algorithm} presents the EM algorithm for Weibull
systems and \Cref{sec:em-results} evaluates its empirical
performance against the direct MLE.
\Cref{sec:scheme1} derives the periodic inspection likelihood
and \Cref{sec:scheme1-results} validates it via Monte Carlo.
\Cref{sec:discussion} discusses practical implications and future
directions.


% ===================================================================
\section{Model}
\label{sec:model}
% ===================================================================

\subsection{Component distributions}

Consider $m$ independent components with lifetimes
$T_1, \ldots, T_m$, where component $j$ follows a Weibull
distribution with shape $\alpha_j > 0$ and scale $\beta_j > 0$:
\begin{equation}\label{eq:weibull-component}
  f_j(t) = \frac{\alpha_j}{\beta_j}
    \Bigl(\frac{t}{\beta_j}\Bigr)^{\alpha_j - 1}
    \exp\!\Bigl[-\Bigl(\frac{t}{\beta_j}\Bigr)^{\alpha_j}\Bigr],
  \qquad
  F_j(t) = 1 - \exp\!\Bigl[-\Bigl(\frac{t}{\beta_j}\Bigr)^{\alpha_j}\Bigr],
\end{equation}
for $t > 0$.  The survival function is
$S_j(t) = \exp[-(t/\beta_j)^{\alpha_j}]$ and the hazard rate is
$h_j(t) = (\alpha_j/\beta_j)(t/\beta_j)^{\alpha_j - 1}$.

The shape parameter $\alpha_j$ governs the failure mechanism:
\begin{itemize}
\item $\alpha_j < 1$: decreasing hazard (early-life failures,
  burn-in);
\item $\alpha_j = 1$: constant hazard (exponential, memoryless);
\item $\alpha_j > 1$: increasing hazard (wear-out, fatigue).
\end{itemize}

The full parameter vector is
$\btheta = (\alpha_1, \beta_1, \ldots, \alpha_m, \beta_m)
\in \Rset_{>0}^{2m}$.

\subsection{\texorpdfstring{$k$-out-of-$n$}{k-out-of-n} systems}

A $k$-out-of-$n$ system functions as long as at least $k$ of its
$n = m$ components survive.  The system lifetime is
\begin{equation}\label{eq:kofn-lifetime}
  T = T_{(n-k+1)},
\end{equation}
where $T_{(1)} \leq \cdots \leq T_{(n)}$ denote the order
statistics of the component lifetimes.  The two extremes are:
\begin{itemize}
\item $k = 1$ (parallel): $T = T_{(n)} = \max_j T_j$;
\item $k = n$ (series): $T = T_{(1)} = \min_j T_j$.
\end{itemize}

The system density for a general coherent system with structure
function $\phi$ is \citep{Samaniego2007}
\begin{equation}\label{eq:sys-density-general}
  f_T(t; \btheta) = \sum_{j=1}^m f_j(t; \alpha_j, \beta_j)
  \sum_{\substack{\mathbf{x}_{-j} : \\ j \text{ critical in }
  (\mathbf{x}_{-j}, x_j = 1)}}
  \prod_{i \neq j} \bigl[F_i(t)\bigr]^{1-x_i}
  \bigl[S_i(t)\bigr]^{x_i},
\end{equation}
where the inner sum ranges over states $\mathbf{x}_{-j}$ of the
other components in which component $j$ is pivotal: flipping $j$
from up to down causes system failure.  For a $k$-out-of-$n$
system, the number of critical states for each component is
$\binom{n-1}{k-1}$ (exactly $k-1$ other components must be up),
and the total number of terms in~\eqref{eq:sys-density-general}
is $m \binom{n-1}{k-1}$.

\subsection{Censoring interpretation}
\label{sec:censoring-interp}

Under Scheme~0 (black-box observation), we observe $n$
independent copies of the system lifetime $T$ with no
component-level information.  As shown in \citet{Towell2026}, the
censoring equivalence for a $k$-out-of-$n$ system yields:
\begin{enumerate}
\item[(a)] One component (the $(n-k+1)$-th to fail) is observed
  exactly: $T_J = T$ for a latent index $J$.
\item[(b)] The $n-k$ components that failed before the system are
  left-censored at $T$: $T_j < T$ for these indices.
\item[(c)] The $k-1$ components still functioning at system failure
  are right-censored at $T$: $T_j > T$ for these indices.
\item[(d)] Component identities within each group are latent.
\end{enumerate}
For parallel systems ($k = 1$): one exact, $m-1$ left-censored,
zero right-censored.  For series systems ($k = n$): one exact,
zero left-censored, $m-1$ right-censored.

This censoring interpretation is distribution-free---it holds for
Weibull components just as it does for exponential ones.  What
changes is the tractability of the resulting likelihood.


% ===================================================================
\section{Likelihood via Critical-State Decomposition}
\label{sec:likelihood}
% ===================================================================

\subsection{System density}

For a $k$-out-of-$n$ system with heterogeneous Weibull components,
the system density~\eqref{eq:sys-density-general} can be written as
\begin{equation}\label{eq:sys-density-kofn}
  f_T(t; \btheta) = \sum_{j=1}^m f_j(t)
  \underbrace{
    \sum_{\substack{S \subseteq \{1,\ldots,m\} \setminus \{j\} \\
    |S| = k-1}}
    \prod_{i \in S} S_i(t) \prod_{i \notin S \cup \{j\}} F_i(t)
  }_{C_j(t; \btheta)},
\end{equation}
where $C_j(t; \btheta)$ is the \emph{critical-state weight} for
component $j$: the probability (given $T = t$) that component $j$
is pivotal, computed over all configurations in which exactly $k-1$
other components survive.

\begin{remark}[Exponential special case]
When $\alpha_j = 1$ for all $j$ (exponential components), each
factor $F_i(t) = 1 - e^{-\lambda_i t}$ and
$S_i(t) = e^{-\lambda_i t}$, and for the parallel case ($k = 1$)
the critical-state weight reduces to
$C_j(t) = \prod_{i \neq j} F_i(t) = \prod_{i \neq j}
(1 - e^{-\lambda_i t})$, which admits the inclusion-exclusion
expansion of \citet{Towell2026}.  For Weibull components with
$\alpha_j \neq 1$, no such expansion is available.
\end{remark}

\subsection{Log-likelihood under Scheme~0}

Given $n$ independent system lifetime observations
$t_1, \ldots, t_n$, the log-likelihood under Scheme~0 is
\begin{equation}\label{eq:loglik-scheme0}
  \ell(\btheta) = \sum_{i=1}^n \ln f_T(t_i; \btheta).
\end{equation}

Each evaluation of $\ell(\btheta)$ requires computing
$f_T(t_i; \btheta)$ via~\eqref{eq:sys-density-kofn}, which
involves:
\begin{itemize}
\item $m$ evaluations of $f_j(t_i)$ (Weibull densities);
\item $m$ evaluations each of $F_j(t_i)$ and $S_j(t_i)$;
\item $m \binom{m-1}{k-1}$ products for the critical-state
  weights.
\end{itemize}
The total cost per observation is
$O\bigl(m^2 \binom{m-1}{k-1}\bigr)$, which is polynomial in $m$
for fixed $k$ but grows as $O(m \cdot 2^{m-1})$ in the worst case
($k \approx m/2$).  For $m = 5$, this is entirely tractable.

\subsection{Gradient computation}

The gradient $\nabla_{\btheta} \ell$ is needed for gradient-based
optimization.  The score with respect to the shape parameter
$\alpha_j$ is
\begin{equation}\label{eq:score-alpha}
  \frac{\partial \ell}{\partial \alpha_j}
  = \sum_{i=1}^n \frac{1}{f_T(t_i)}
    \frac{\partial f_T(t_i)}{\partial \alpha_j},
\end{equation}
where $\partial f_T / \partial \alpha_j$ propagates through both
$f_j(t)$ and the CDF/survival terms $F_i(t)$, $S_i(t)$ that appear
in the critical-state weights of \emph{every} component $k$ (since
$\alpha_j$ affects $F_j$ and $S_j$, which appear in $C_k$ for
$k \neq j$).

The partial derivatives of the Weibull density and CDF with respect
to the shape parameter are:
\begin{align}
  \frac{\partial f_j(t)}{\partial \alpha_j}
  &= f_j(t) \Bigl[
    \frac{1}{\alpha_j}
    + \ln\!\Bigl(\frac{t}{\beta_j}\Bigr)
    - \Bigl(\frac{t}{\beta_j}\Bigr)^{\alpha_j}
      \ln\!\Bigl(\frac{t}{\beta_j}\Bigr)
  \Bigr],
  \label{eq:df-dalpha} \\
  \frac{\partial F_j(t)}{\partial \alpha_j}
  &= S_j(t) \Bigl(\frac{t}{\beta_j}\Bigr)^{\alpha_j}
    \ln\!\Bigl(\frac{t}{\beta_j}\Bigr).
  \label{eq:dF-dalpha}
\end{align}
Analogous expressions hold for $\partial / \partial \beta_j$.
These are elementary functions of $t$ but involve $\ln(t/\beta_j)$
terms that prevent the exponential cancellations that make the
exponential case tractable.

In practice, numerical differentiation via finite differences
(Richardson extrapolation) provides a robust alternative to
analytical gradients, at the cost of additional likelihood
evaluations.  Our implementation uses the \texttt{numDeriv} package
in R for Hessian computation and the L-BFGS-B algorithm
\citep{ByrdLuNocedal1995} with box constraints
$\alpha_j, \beta_j > 0$ for optimization.


% ===================================================================
\section{Why Exponential Results Don't Generalize}
\label{sec:why-hard}
% ===================================================================

The difficulties with Weibull system estimation are not merely
computational overhead atop an exponential foundation.  They
represent structural changes in the mathematical problem that
invalidate the entire toolkit developed for the exponential case.

\subsection{The inclusion-exclusion trick fails}

For exponential components, the parallel system CDF factors as
\begin{equation}\label{eq:exp-cdf-ie}
  F_T(t) = \prod_{j=1}^m (1 - e^{-\lambda_j t})
  = \sum_{S \subseteq \{1,\ldots,m\}} (-1)^{|S|}
    e^{-\lambda(S)\, t},
\end{equation}
where $\lambda(S) = \sum_{j \in S} \lambda_j$.  Each term is a
simple exponential, and integrals of the form
$\int_a^b e^{-\lambda(S)\, t}\, dt = (e^{-\lambda(S)\, a} -
e^{-\lambda(S)\, b}) / \lambda(S)$ are elementary.  This is the
technical engine behind the closed-form likelihoods, the
analytical EM, and the identifiability proof of \citet{Towell2026}.

For Weibull components, the analogous product is
\begin{equation}\label{eq:weibull-cdf-product}
  F_T(t) = \prod_{j=1}^m
    \Bigl(1 - e^{-(t/\beta_j)^{\alpha_j}}\Bigr).
\end{equation}
Expanding this product via inclusion-exclusion yields terms of the
form
\[
  (-1)^{|S|} \exp\!\Bigl[-\sum_{j \in S}
    \Bigl(\frac{t}{\beta_j}\Bigr)^{\alpha_j}\Bigr].
\]
When the $\alpha_j$ are not all equal, the exponent is a sum of
different power functions of $t$---not a linear function of $t$.
The integral $\int_a^b \exp[-\sum_{j \in S}
(t/\beta_j)^{\alpha_j}]\, dt$ has no closed form in general.  In
the special case $\alpha_j = \alpha$ for all $j$ (common shape),
the exponent becomes $-t^\alpha \sum_{j \in S} \beta_j^{-\alpha}$,
and a substitution $u = t^\alpha$ reduces the integral to an
incomplete gamma function---tractable but not elementary.

\subsection{Series hazard additivity fails}

For exponential series systems:
\[
  h_T(t) = \sum_{j=1}^m \lambda_j = \Lambda \quad
  \text{(constant)}.
\]
The system is exponential with known rate $\Lambda$; individual
rates are unidentifiable, but the aggregate is trivially estimated.

For Weibull series systems:
\[
  h_T(t) = \sum_{j=1}^m \frac{\alpha_j}{\beta_j}
    \Bigl(\frac{t}{\beta_j}\Bigr)^{\alpha_j - 1}.
\]
This is a sum of power functions with different exponents.  The
system is no longer Weibull (unless all $\alpha_j$ are equal).
Individual $(\alpha_j, \beta_j)$ pairs are not identifiable from
$T_{\min}$ alone, and the mixture of power-law hazards creates a
likelihood surface with multiple local optima.  Our experiments
confirm catastrophic behavior: RMSE exceeding 360 for mixed-shape
series systems (\Cref{sec:sim-results}).

\subsection{The memoryless property fails}

The exponential distribution's memoryless property---$P(T > s + t
\mid T > s) = P(T > t)$---implies that the conditional
distribution of an exponential component given survival to time $s$
is again exponential with the same rate.  This is what makes the EM
algorithm's E-step tractable: the truncated exponential mean
$\E[T_j \mid T_j < t] = 1/\lambda_j -
t/(e^{\lambda_j t} - 1)$ has a simple closed form.

For Weibull components, the conditional distribution given
$T_j < t$ has density
\begin{equation}\label{eq:truncated-weibull}
  g(u \mid T_j < t) = \frac{f_j(u)}{F_j(t)}, \qquad 0 < u < t.
\end{equation}
The conditional mean is
\begin{equation}\label{eq:truncated-weibull-mean}
  \E[T_j \mid T_j < t]
  = \frac{\int_0^t u\, f_j(u)\, du}{F_j(t)}
  = \frac{\beta_j\, \gamma(1 + 1/\alpha_j,\;
    (t/\beta_j)^{\alpha_j})}{F_j(t)\,
    \Gamma(1 + 1/\alpha_j)}\; \cdot\;
  \E[T_j],
\end{equation}
where $\gamma(a, x) = \int_0^x u^{a-1} e^{-u}\, du$ is the lower
incomplete gamma function.  This is computationally available
(standard libraries implement $\gamma$) but not elementary, and it
depends on both $\alpha_j$ and $\beta_j$ in a way that prevents
the simple reciprocal M-step of the exponential EM.

\subsection{The identifiability landscape changes}

For exponential parallel systems, \citet{BasuGhosh1980} and
\citet{Towell2026} proved identifiability: the distribution of
$\max_j T_j$ uniquely determines the component rates.  The proof
extracts rates one at a time from the tail behavior of the reversed
hazard rate.

For Weibull parallel systems, identifiability of the full
parameter vector $\btheta$ is a more delicate question.
Consider two components with parameters $(\alpha_1, \beta_1)$ and
$(\alpha_2, \beta_2)$.  The parallel system CDF is
\[
  F_T(t) = (1 - e^{-(t/\beta_1)^{\alpha_1}})
            (1 - e^{-(t/\beta_2)^{\alpha_2}}).
\]
If $\alpha_1 = \alpha_2 = \alpha$ (common shape), then the same
tail-extraction argument applies: the reversed hazard rate
$r_T(t) = f_T(t)/F_T(t)$ is asymptotically dominated by the
component with the largest scale (smallest hazard at large $t$),
and $\alpha$ and the largest $\beta_j$ can be extracted.  For
heterogeneous shapes, the tail behavior involves competing power
laws, and the extraction is more complex.  A full identifiability
proof for the heterogeneous Weibull case is beyond the scope of this
paper; we note only that the numerical evidence (converged MLEs with
reasonable accuracy when $k$ is chosen appropriately) is consistent
with identifiability for the configurations tested.


% ===================================================================
\section{Direct MLE: Simulation Study}
\label{sec:simulations}
% ===================================================================

We evaluate the direct MLE for Weibull components in $k$-out-of-$n$
systems through a systematic Monte Carlo study.  The goal is to
quantify the degradation in estimation quality relative to the
exponential baseline and to identify the dependence of optimal
system structure on the shape parameter.

\subsection{Experimental design}
\label{sec:sim-design}

\paragraph{Components.}
We use $m = 5$ components with Weibull parameters
median-matched to the exponential rates
$\blam = (0.5, 0.4, 0.3, 0.2, 0.1)$ used in the companion
exponential study \citep{Towell2026}.  For
$\mathrm{Exp}(\lambda_j)$, the median is
$\log(2)/\lambda_j$; for
$\mathrm{Weibull}(\alpha_j, \beta_j)$, the median is
$\beta_j [\log(2)]^{1/\alpha_j}$.  Setting these equal gives the
median-matching scale:
\begin{equation}\label{eq:median-match}
  \beta_j = \frac{\log(2)/\lambda_j}{[\log(2)]^{1/\alpha_j}}.
\end{equation}

\paragraph{Shape scenarios.}
Three scenarios capture qualitatively different failure physics:
\begin{enumerate}
\item \textbf{Exponential baseline} ($\alpha_j = 1$ for all $j$):
  validates the framework against known exponential results.
\item \textbf{Increasing hazard} ($\alpha_j = 2$ for all $j$):
  wear-out components with identical shapes but different scales.
\item \textbf{Mixed shapes}
  ($\balpha = (2.0, 1.5, 1.0, 0.8, 0.5)$): the most realistic
  scenario, combining wear-out, constant-hazard, and early-life
  components.
\end{enumerate}

\paragraph{System structures.}
For each scenario, we test all five $k$-out-of-5 structures:
$k = 1$ (parallel) through $k = 5$ (series).

\paragraph{Estimation.}
Each cell uses $n = 200$ system lifetime observations and
$R = 15$ Monte Carlo replications.  The MLE is computed by
L-BFGS-B with box constraints $\alpha_j, \beta_j > 10^{-10}$,
with Nelder--Mead \citep{NelderMead1965} fallback on a
log-transformed parameterization.  Multi-start optimization
($n_{\text{starts}} = 5$) with random perturbations of the initial
values guards against local optima.  Standard errors are obtained
from the observed Fisher information via Richardson-extrapolated
numerical Hessian.

\paragraph{Initialization.}
Starting values are $\alpha_j^{(0)} = 1$ (agnostic shape) and
$\beta_j^{(0)} = \bar{t} \cdot c_j$ where $c_j$ is a spread
factor $c_j = 0.5 + (j-1)/(m-1)$, $j = 1, \ldots, m$.  The
spread breaks the permutation symmetry of the likelihood under
Scheme~0, analogous to the exponential case \citep{Towell2026}.

\paragraph{Metrics.}
We report:
\begin{itemize}
\item Convergence rate (\% of replications where the optimizer
  reports convergence with finite parameter estimates);
\item Mean RMSE over scale parameters:
  $m^{-1} \sum_j \bigl[R_c^{-1}
  \sum_r (\hat\beta_{j}^{(r)} - \beta_j)^2\bigr]^{1/2}$,
  where $R_c$ is the number of converged replications;
\item Maximum RMSE over scale parameters:
  $\max_j \bigl[R_c^{-1}
  \sum_r (\hat\beta_{j}^{(r)} - \beta_j)^2\bigr]^{1/2}$.
\end{itemize}
Scale RMSE is the primary metric because scale parameters are on
the same physical scale as the component lifetimes.  Shape
parameter RMSE is reported separately where relevant.

\paragraph{Label matching.}
Since the Scheme~0 likelihood is invariant to permutation of the
component labels, estimated parameters are sorted by decreasing
scale and compared to the correspondingly sorted true values.

\subsection{Results}
\label{sec:sim-results}

\Cref{tab:weibull-results} presents the full results.
Three patterns dominate.

\begin{table}[t]
\centering
\caption{%
  Weibull $k$-out-of-5 estimation: convergence rate, mean scale RMSE,
  and maximum scale RMSE across three shape scenarios
  ($n = 200$, $R = 15$ replications per cell).  The optimal $k$ for
  each scenario (by mean RMSE, excluding catastrophic cases) is marked
  with $\star$.
}
\label{tab:weibull-results}
\small
\begin{tabular}{cl rrr}
\toprule
$k$ & Scenario & Conv.\ (\%) & Mean RMSE & Max RMSE \\
\midrule
\multirow{3}{*}{1}
  & Baseline ($\alpha = 1$)
    & 100 & 1.846 & 2.500 \\
  & Increasing ($\alpha = 2$)
    & 93 & \textbf{1.113}$\,\star$ & 1.586 \\
  & Mixed shapes
    & 73 & 2.639 & 5.175 \\
\addlinespace
\multirow{3}{*}{2}
  & Baseline ($\alpha = 1$)
    & 80 & 1.965 & 3.327 \\
  & Increasing ($\alpha = 2$)
    & 80 & 1.248 & 2.544 \\
  & Mixed shapes
    & 73 & \textbf{1.994}$\,\star$ & 5.350 \\
\addlinespace
\multirow{3}{*}{3}
  & Baseline ($\alpha = 1$)
    & 73 & 1.795 & 3.843 \\
  & Increasing ($\alpha = 2$)
    & 60 & 1.385 & 3.789 \\
  & Mixed shapes
    & 80 & 2.564 & 8.632 \\
\addlinespace
\multirow{3}{*}{4}
  & Baseline ($\alpha = 1$)
    & 67 & \textbf{1.760}$\,\star$ & 3.956 \\
  & Increasing ($\alpha = 2$)
    & 53 & 7.497 & 12.568 \\
  & Mixed shapes
    & 47 & 8.489 & 16.163 \\
\addlinespace
\multirow{3}{*}{5}
  & Baseline ($\alpha = 1$)
    & 93 & 23.717 & 95.803 \\
  & Increasing ($\alpha = 2$)
    & 80 & 1.554 & 5.185 \\
  & Mixed shapes
    & 73 & 363.438 & 1540.711 \\
\bottomrule
\end{tabular}
\end{table}

\subsubsection{Finding 1: Weibull estimation is 10--30 times harder
  than exponential}
\label{sec:finding1}

The exponential baseline with $m = 5$ and $n = 200$ achieves
RMSE $\approx 0.07\text{--}0.14$ across $k$ in
\citet{Towell2026}, with 100\% convergence throughout.  The
Weibull baseline (same rates, same $n$, but 10 parameters instead
of 5) has best RMSE $= 1.760$ (at $k = 4$), a degradation factor
exceeding 12.  Convergence drops to 67--100\% depending on $k$.

The factor-of-two increase in parameter count (from $m$ to $2m$)
does not explain a 12-fold RMSE increase.  The additional
difficulty comes from the likelihood surface geometry:
shape--scale confounding creates ridges and local optima that
numerical optimizers struggle to navigate.

\subsubsection{Finding 2: Shape parameter changes the optimal \texorpdfstring{$k$}{k}}
\label{sec:sim-optimal-k}

The optimal system structure for estimation depends on the shape
parameter:
\begin{center}
\begin{tabular}{lcc}
\toprule
Scenario & Optimal $k^*$ & Best RMSE \\
\midrule
Baseline ($\alpha = 1$) & 4 & 1.760 \\
Increasing ($\alpha = 2$) & 1 & 1.113 \\
Mixed shapes & 2 & 1.994 \\
\bottomrule
\end{tabular}
\end{center}

This is a central finding: \emph{a practitioner who designs a test
system based on exponential theory would choose $k = 4$, but the
correct choice for wear-out components is $k = 1$ (parallel)}.
The intuition is that increasing-hazard components fail in a more
predictable temporal sequence---the component with the largest scale
tends to fail last with high probability---so the parallel system
(which observes the maximum) already extracts substantial ordering
information.  Under constant hazard, the failure sequence is more
random, and higher $k$ (which reveals more of the failure ordering)
provides additional information.

\subsubsection{Finding 3: Series systems are catastrophic}

The series case ($k = 5$) shows:
\begin{itemize}
\item Baseline: RMSE $= 23.717$ (13 times the best);
\item Increasing hazard: RMSE $= 1.554$ (close to best, but with
  only 80\% convergence);
\item Mixed shapes: RMSE $= 363.438$ (182 times the best).
\end{itemize}

The catastrophe under mixed shapes is particularly striking.  The
five shape parameters create a 10-dimensional likelihood with
extensive multi-modality when only $T_{\min}$ is observed.  The
optimizer converges to wildly different local optima across
replications, inflating RMSE by three orders of magnitude.

The increasing-hazard series case is an exception: with common
shape $\alpha = 2$, the system hazard
$h_T(t) = \sum_j (2/\beta_j)(t/\beta_j)$ is a linear function of
$t$ (since all shapes are 2), and the problem is structurally
simpler.  But even here, convergence (80\%) is below the parallel
case (93\%).

\subsubsection{Finding 4: Convergence degrades systematically}

For the increasing-hazard scenario, convergence decreases
monotonically with $k$: $93\% \to 80\% \to 60\% \to 53\%$
(for $k = 1, 2, 3, 4$).  Higher $k$ means the system fails later
(more components must fail), producing longer-tailed observations
where the Weibull shape parameter has its strongest effect and the
likelihood surface is sharpest.  The $k = 5$ (series) case
partially recovers to 80\% because the common-shape structure
simplifies the optimization landscape.

For mixed shapes, the worst convergence is 47\% at $k = 4$.  Fewer
than half the replications produce a usable estimate---a level of
unreliability that rules out practical application without
algorithmic improvements.

\subsection{Comparison with exponential baseline}

\Cref{tab:exp-comparison} compares the Weibull and exponential
cases directly.

\begin{table}[t]
\centering
\caption{%
  Direct comparison: exponential ($m = 5$ parameters) vs.\ Weibull
  baseline ($2m = 10$ parameters, $\alpha = 1$) at the same
  median-matched rates and $n = 200$.  Exponential results from
  \citet{Towell2026}.
}
\label{tab:exp-comparison}
\small
\begin{tabular}{c rr rr}
\toprule
& \multicolumn{2}{c}{Exponential} &
  \multicolumn{2}{c}{Weibull ($\alpha = 1$)} \\
\cmidrule(lr){2-3} \cmidrule(lr){4-5}
$k$ & Conv.\ (\%) & RMSE & Conv.\ (\%) & RMSE \\
\midrule
1 & 100 & $\sim$0.09 & 100 & 1.846 \\
2 & 100 & $\sim$0.11 & 80  & 1.965 \\
3 & 100 & $\sim$0.07 & 73  & 1.795 \\
4 & 100 & $\sim$0.14 & 67  & 1.760 \\
5 & 100 & $\sim$0.08 & 93  & 23.717 \\
\bottomrule
\end{tabular}
\end{table}

The comparison is stark: the ``same'' estimation problem---with
identical rates but modeled as Weibull (shape~$= 1$)---is
20 times worse in RMSE and 7--33 percentage points worse in
convergence.  The degradation comes entirely from the additional
shape parameters, which the optimizer must estimate even though
their true values are~1.  This illustrates the cost of model
flexibility: even when the true model is exponential, fitting a
Weibull model pays a substantial penalty in finite samples.


% ===================================================================
\section{EM Algorithm for Weibull Systems}
\label{sec:em-algorithm}
% ===================================================================

The convergence difficulties of direct MLE for Weibull systems
motivate an EM algorithm that exploits the censoring structure of
the $k$-out-of-$n$ system.  The EM is attractive because it
guarantees monotone likelihood increase at every iteration
\citep{DempsterLairdRubin1977, Wu1983}, replacing a single hard
optimization with a sequence of easier sub-problems.

\subsection{Complete and observed data}

For a $k$-out-of-$n$ system with Weibull components, the complete
data for one system observation consists of:
\begin{itemize}
\item The component lifetimes $(T_1, \ldots, T_m)$;
\item The system lifetime $T = T_{(n-k+1)}$;
\item The identity of the critical component
  $J = \argmax\{j : T_j = T_{(n-k+1)}\}$;
\item The group assignments: which components are left-censored
  ($T_j < T$), which is exact ($T_j = T$), and which are
  right-censored ($T_j > T$).
\end{itemize}
The observed data under Scheme~0 is only $T$; all component-level
information is latent.

\subsection{Complete-data log-likelihood}

Given complete data $(T_1, \ldots, T_m)$ for $n$ independent
system observations, the complete-data log-likelihood is
\begin{equation}\label{eq:complete-loglik}
  \ell_c(\btheta) = \sum_{i=1}^n \sum_{j=1}^m
    \ln f_j(T_{ij}; \alpha_j, \beta_j)
  = \sum_{j=1}^m \sum_{i=1}^n \Bigl[
    \ln \frac{\alpha_j}{\beta_j}
    + (\alpha_j - 1) \ln \frac{T_{ij}}{\beta_j}
    - \Bigl(\frac{T_{ij}}{\beta_j}\Bigr)^{\alpha_j}
  \Bigr],
\end{equation}
which decomposes into $m$ independent Weibull log-likelihoods, one
per component.  This separability is the key advantage of the EM:
the M-step reduces to $m$ independent Weibull MLEs.

\subsection{E-step: conditional expectations}

Given observed system lifetime $T = t$ and current parameter
estimates $\btheta^{(s)}$, the E-step computes the expected
complete-data sufficient statistics.  For the Weibull family, the
sufficient statistics for component $j$ are not simply the
lifetimes $T_{ij}$ but the functions $\ln T_{ij}$ and
$T_{ij}^{\alpha_j}$ (since the Weibull is an exponential family
in the variable $T^{\alpha}$ with natural parameter
$-\beta^{-\alpha}$).  However, the M-step does not have a
closed-form solution even for complete Weibull data---it requires
solving a transcendental equation for $\alpha_j$---so the E-step
need only provide the conditional distribution (or sufficient
functionals thereof) for each component lifetime.

\subsubsection{Classification probabilities}

For a $k$-out-of-$n$ system, the system observation $T = t$
partitions components into three groups:
\begin{enumerate}
\item The critical component ($T_j = t$);
\item Left-censored components ($T_j < t$);
\item Right-censored components ($T_j > t$).
\end{enumerate}
The probability that component $j$ is critical, given $T = t$, is:
\begin{equation}\label{eq:em-pij}
  p_j(t; \btheta^{(s)})
  = \frac{f_j(t; \btheta^{(s)}) \cdot C_j(t; \btheta^{(s)})}
         {f_T(t; \btheta^{(s)})},
\end{equation}
where $C_j$ is the critical-state weight
from~\eqref{eq:sys-density-kofn}.  These probabilities satisfy
$\sum_j p_j(t; \btheta^{(s)}) = 1$.

\subsubsection{Conditional moments for the E-step}

For the M-step, we need the following conditional expectations for
each observation $i$ and component $j$:

\paragraph{Case 1: Component $j$ is critical ($T_j = t_i$).}
\begin{equation}\label{eq:em-exact-case}
  \E[\ln T_j \mid T_j = t_i] = \ln t_i, \qquad
  \E[T_j^{\alpha_j} \mid T_j = t_i] = t_i^{\alpha_j}.
\end{equation}

\paragraph{Case 2: Component $j$ is left-censored ($T_j < t_i$).}
The conditional distribution is
$f_j(u) / F_j(t_i)$ on $(0, t_i)$.  The required moments are:
\begin{align}
  \E[\ln T_j \mid T_j < t_i; \btheta^{(s)}]
  &= \frac{\int_0^{t_i} (\ln u)\, f_j(u)\, du}{F_j(t_i)},
  \label{eq:em-left-lnT} \\
  \E[T_j^{\alpha_j} \mid T_j < t_i; \btheta^{(s)}]
  &= \frac{\int_0^{t_i} u^{\alpha_j}\, f_j(u)\, du}{F_j(t_i)}.
  \label{eq:em-left-Talpha}
\end{align}

The integral in~\eqref{eq:em-left-Talpha} can be expressed in
terms of the incomplete gamma function.  With the substitution
$v = (u/\beta_j)^{\alpha_j}$:
\begin{equation}\label{eq:em-left-Talpha-gamma}
  \int_0^{t_i} u^{\alpha_j}\, f_j(u)\, du
  = \beta_j^{\alpha_j}
    \int_0^{(t_i/\beta_j)^{\alpha_j}}
    v\, e^{-v}\, dv
  = \beta_j^{\alpha_j}\, \gamma\!\bigl(2,\;
    (t_i/\beta_j)^{\alpha_j}\bigr),
\end{equation}
where $\gamma(a, x) = \int_0^x u^{a-1} e^{-u}\, du$ is the
lower incomplete gamma function.

The integral in~\eqref{eq:em-left-lnT} is more complex but can be
expressed as:
\begin{equation}\label{eq:em-left-lnT-formula}
  \int_0^{t_i} (\ln u)\, f_j(u)\, du
  = \frac{\ln \beta_j}{\alpha_j}\, \gamma\!\bigl(1,\;
    (t_i/\beta_j)^{\alpha_j}\bigr)
  + \frac{1}{\alpha_j}
    \int_0^{(t_i/\beta_j)^{\alpha_j}}
    (\ln v)\, e^{-v}\, dv.
\end{equation}
The remaining integral involves
$\int_0^x (\ln v) e^{-v}\, dv$, which is related to the
derivative of the incomplete gamma function with respect to its
first argument.  In practice, this is most reliably computed by
numerical quadrature.

\paragraph{Case 3: Component $j$ is right-censored ($T_j > t_i$).}
The conditional distribution is
$f_j(u) / S_j(t_i)$ on $(t_i, \infty)$.  The required
expectations are:
\begin{align}
  \E[\ln T_j \mid T_j > t_i; \btheta^{(s)}]
  &= \frac{\int_{t_i}^\infty (\ln u)\, f_j(u)\, du}{S_j(t_i)},
  \label{eq:em-right-lnT} \\
  \E[T_j^{\alpha_j} \mid T_j > t_i; \btheta^{(s)}]
  &= \frac{\beta_j^{\alpha_j}\,
    \Gamma\!\bigl(2,\; (t_i/\beta_j)^{\alpha_j}\bigr)}
    {S_j(t_i)},
  \label{eq:em-right-Talpha}
\end{align}
where $\Gamma(a, x) = \int_x^\infty u^{a-1} e^{-u}\, du$ is the
upper incomplete gamma function.

\subsubsection{Marginalization over group assignments}

For a general $k$-out-of-$n$ system, component $j$ may be in any
of the three groups (critical, left-censored, right-censored),
depending on the latent assignment.  The E-step marginalizes over
all possible assignments.  For observation $i$ with system lifetime
$t_i$:

\begin{equation}\label{eq:em-marginalized}
  \E[g(T_j) \mid T = t_i; \btheta^{(s)}]
  = p_j \cdot g(t_i)
  + q_j^L \cdot \E[g(T_j) \mid T_j < t_i]
  + q_j^R \cdot \E[g(T_j) \mid T_j > t_i],
\end{equation}
where $p_j$ is the probability that $j$ is critical, $q_j^L$ is
the probability that $j$ is left-censored, and $q_j^R$ is the
probability that $j$ is right-censored, all conditional on
$T = t_i$ and current parameters $\btheta^{(s)}$.  For a parallel
system ($k = 1$), $q_j^R = 0$ for all $j$; for a series system
($k = n$), $q_j^L = 0$ for all $j$.

The classification probabilities $(p_j, q_j^L, q_j^R)$ are
computed from the critical-state decomposition:
$p_j$ is as in~\eqref{eq:em-pij}, and $q_j^L + q_j^R = 1 - p_j$.
To split $1 - p_j$ between left- and right-censoring, we compute
the probability that component $j$ has failed by time $t$
\emph{and} is not critical:
\begin{equation}\label{eq:em-qjL}
  q_j^L = \frac{F_j(t) \cdot D_j^L(t; \btheta^{(s)})}
               {f_T(t; \btheta^{(s)})},
  \qquad
  q_j^R = 1 - p_j - q_j^L,
\end{equation}
where $D_j^L(t; \btheta^{(s)})$ sums over critical states in
which component $j$ is down and not pivotal.

\subsection{M-step: component-wise Weibull MLE}

Given the expected sufficient statistics from the E-step, the
M-step solves $m$ independent Weibull maximum likelihood problems.
For component $j$, the expected complete-data log-likelihood
(dropping constants) is:
\begin{equation}\label{eq:em-mstep-obj}
  Q_j(\alpha_j, \beta_j)
  = n \ln \frac{\alpha_j}{\beta_j}
  + (\alpha_j - 1) \sum_{i=1}^n \hat\mu_{ij}^{(\ln)}
  - \frac{1}{\beta_j^{\alpha_j}} \sum_{i=1}^n \hat\mu_{ij}^{(\alpha)},
\end{equation}
where
$\hat\mu_{ij}^{(\ln)} = \E[\ln T_j \mid T = t_i; \btheta^{(s)}]$
and
$\hat\mu_{ij}^{(\alpha)} = \E[T_j^{\alpha_j^{(s)}} \mid T = t_i;
\btheta^{(s)}]$
are from the E-step.

Setting $\partial Q_j / \partial \beta_j = 0$ gives:
\begin{equation}\label{eq:em-beta-update}
  \beta_j^{(s+1)} = \Bigl(
    \frac{1}{n} \sum_{i=1}^n \hat\mu_{ij}^{(\alpha)}
  \Bigr)^{1/\alpha_j^{(s+1)}}.
\end{equation}
Setting $\partial Q_j / \partial \alpha_j = 0$ gives the
profile equation:
\begin{equation}\label{eq:em-alpha-profile}
  \frac{1}{\alpha_j^{(s+1)}}
  + \frac{1}{n} \sum_{i=1}^n \hat\mu_{ij}^{(\ln)}
  = \frac{\sum_{i=1}^n \hat\mu_{ij}^{(\alpha)}
    \cdot \hat\nu_{ij}^{(\ln)}}
   {\sum_{i=1}^n \hat\mu_{ij}^{(\alpha)}},
\end{equation}
where
$\hat\nu_{ij}^{(\ln)} = \E[T_j^{\alpha_j} \ln T_j \mid T = t_i;
\btheta^{(s)}]$.  This is a nonlinear equation in $\alpha_j$ that
must be solved numerically (e.g., by Newton--Raphson or bisection),
as is standard for single-component Weibull MLE
\citep{Cohen1965, LawlessBook2003}.

\begin{remark}[Contrast with exponential EM]
For exponential components, the E-step involves the truncated
exponential mean $1/\lambda_j - t/(e^{\lambda_j t} - 1)$ (a
simple closed form), and the M-step is $\lambda_j^{(s+1)} = n /
\sum_i \hat\tau_{ij}$ (a reciprocal).  For Weibull components, the
E-step requires incomplete gamma functions (which are available
in standard numerical libraries via the regularized gamma function),
and the M-step requires a one-dimensional numerical optimization
over $\alpha_j$ with closed-form $\beta_j$ given $\alpha_j$
via~\eqref{eq:em-beta-update}.  The EM retains its structural
advantage (monotone likelihood increase, decomposition into $m$
sub-problems) at the cost of replacing elementary closed forms
with well-conditioned numerical computations.
\Cref{sec:em-results} evaluates this algorithm empirically.
\end{remark}

\subsection{Algorithm summary}

\Cref{alg:weibull-em} provides pseudocode for the complete EM
algorithm.

\begin{algorithm}[t]
\caption{EM algorithm for Weibull $k$-out-of-$n$ system estimation}
\label{alg:weibull-em}
\begin{algorithmic}[1]
\Require System observations $t_1, \ldots, t_n > 0$;\;
  system structure $(k, n)$;\;
  tolerance $\varepsilon > 0$
\Ensure MLE $\hat\btheta =
  (\hat\alpha_1, \hat\beta_1, \ldots, \hat\alpha_m, \hat\beta_m)$
\State Initialize $\alpha_j^{(0)} \gets 1$,\;
  $\beta_j^{(0)} \gets \bar{t} \cdot (0.5 + (j-1)/(m-1))$
  for $j = 1, \ldots, m$
\State $s \gets 0$
\Repeat
  \For{$i = 1, \ldots, n$}
    \Comment{\textbf{E-step}}
    \State Compute $f_T(t_i; \btheta^{(s)})$ via
      critical-state decomposition~\eqref{eq:sys-density-kofn}
    \For{$j = 1, \ldots, m$}
      \State Compute classification probabilities
        $(p_j, q_j^L, q_j^R)$
        via~\eqref{eq:em-pij}--\eqref{eq:em-qjL}
      \State Compute
        $\hat\mu_{ij}^{(\ln)}$,
        $\hat\mu_{ij}^{(\alpha)}$
        via~\eqref{eq:em-marginalized}
        using~\eqref{eq:em-exact-case}--\eqref{eq:em-right-Talpha}
    \EndFor
  \EndFor
  \For{$j = 1, \ldots, m$}
    \Comment{\textbf{M-step}}
    \State Solve~\eqref{eq:em-alpha-profile} for
      $\alpha_j^{(s+1)}$ (Newton--Raphson)
    \State Compute $\beta_j^{(s+1)}$
      via~\eqref{eq:em-beta-update}
  \EndFor
  \State $s \gets s + 1$
\Until{$|\ell(\btheta^{(s)}) - \ell(\btheta^{(s-1)})| /
  |\ell(\btheta^{(s-1)})| < \varepsilon$}
\State \Return $\hat\btheta \gets \btheta^{(s)}$
\end{algorithmic}
\end{algorithm}


% ===================================================================
\section{Periodic Inspection: Scheme 1}
\label{sec:scheme1}
% ===================================================================

The simulation results of \Cref{sec:simulations} demonstrate that
Scheme~0 (system-level-only observation) is often insufficient for
reliable Weibull parameter recovery.  A natural remedy is to
supplement system failure times with component-level information
obtained through periodic inspection.

\subsection{Observation model}

Under Scheme~1, components are inspected at regular intervals
$\tau, 2\tau, 3\tau, \ldots$ while the system is operating.  At
each inspection time $\ell\tau$ (with $\ell = 1, 2, \ldots$), the
status of each component (functioning or failed) is recorded.  The
system is observed until failure, at which point the exact system
failure time $T$ is recorded.

For each system $i$ with failure time $T_i$, and for each
component $j$, the inspection data determines an interval
$(\underline{t}_{ij}, \bar{t}_{ij}]$ containing the component
failure time:
\begin{itemize}
\item $\underline{t}_{ij} = \ell\tau$ if component $j$ was
  functioning at the last inspection before system failure and
  failed by time $T_i$;
\item $\bar{t}_{ij} = (\ell + 1)\tau$ or $T_i$, whichever is
  smaller;
\item If component $j$ is still functioning at system failure
  ($T_j > T_i$), the observation is right-censored at $T_i$.
\end{itemize}

\subsection{Scheme~1 likelihood}

The Scheme~1 likelihood for one system observation combines the
exact system failure time with the interval-censored component data.
Let $\mathcal{L}_j$ denote the set of components for which we have
interval-censored data $T_j \in (\underline{t}_{ij},
\bar{t}_{ij}]$, $\mathcal{E}_j$ the component (if any) known to
have failed at $T_i$, and $\mathcal{R}_j$ the components
right-censored at $T_i$.

The likelihood contribution of system $i$ is:
\begin{equation}\label{eq:scheme1-lik}
  L_i(\btheta) = f_T(T_i; \btheta)
  \cdot \prod_{j \in \mathcal{L}_j}
    \frac{F_j(\bar{t}_{ij}) - F_j(\underline{t}_{ij})}
         {\tilde{F}_j(T_i)}
  \cdot \prod_{j \in \mathcal{R}_j}
    \frac{S_j(T_i)}{\tilde{S}_j(T_i)},
\end{equation}
where $\tilde{F}_j$ and $\tilde{S}_j$ are the appropriate
normalizing functions that account for the conditioning on the
system failure time.

In practice, the joint likelihood is most cleanly expressed by
treating the complete observation as:
\begin{equation}\label{eq:scheme1-loglik}
  \ell_1(\btheta) = \sum_{i=1}^n \Bigl[
    \ln f_T(T_i; \btheta)
    + \sum_{j=1}^m \ln P(D_{ij} \mid T_i, \btheta)
  \Bigr],
\end{equation}
where $D_{ij}$ denotes the inspection data for component $j$ in
system $i$, and $P(D_{ij} \mid T_i, \btheta)$ is:
\begin{equation}\label{eq:scheme1-component}
  P(D_{ij} \mid T_i, \btheta) = \begin{cases}
    \dfrac{F_j(\bar{t}_{ij}) - F_j(\underline{t}_{ij})}
          {F_j(T_i)}
    & \text{if } T_j \leq T_i \text{ (interval-censored)}, \\[12pt]
    \dfrac{S_j(T_i)}{S_j(T_i)} = 1
    & \text{if } T_j > T_i \text{ (right-censored)}.
  \end{cases}
\end{equation}

The key observation is that the second term
in~\eqref{eq:scheme1-loglik} provides \emph{component-level}
information that is entirely absent under Scheme~0.  For Weibull
components, the interval $(\underline{t}_{ij}, \bar{t}_{ij}]$
constrains both $\alpha_j$ and $\beta_j$: the probability of
falling in a specific interval depends on the shape (which
determines how the hazard rate changes over the interval) and
the scale (which determines the overall failure timing).  This
joint constraint resolves the shape--scale ambiguity that makes
Scheme~0 estimation so difficult.

\subsection{Information gain from inspection}

The inspection interval width $\tau$ controls the information
content:
\begin{itemize}
\item $\tau \to 0$: continuous monitoring (Scheme~2), complete data,
  trivial estimation;
\item $\tau \to \infty$: no inspections (Scheme~0), minimal
  information;
\item Finite $\tau$: intermediate, with finer intervals
  providing more precise component-level data at higher
  monitoring cost.
\end{itemize}

For the EM algorithm of \Cref{sec:em-algorithm}, the Scheme~1
E-step replaces the unconditional truncated moments with moments
conditional on the interval constraint:
\begin{equation}\label{eq:scheme1-em-estep}
  \E[g(T_j) \mid T_j \in (\underline{t}_{ij}, \bar{t}_{ij}];
  \btheta^{(s)}]
  = \frac{\int_{\underline{t}_{ij}}^{\bar{t}_{ij}}
    g(u)\, f_j(u; \btheta^{(s)})\, du}
    {F_j(\bar{t}_{ij}; \btheta^{(s)}) -
     F_j(\underline{t}_{ij}; \btheta^{(s)})}.
\end{equation}
The interval constraint concentrates the conditional distribution,
leading to more informative expected sufficient statistics and
faster EM convergence.


% ===================================================================
\section{Empirical Results: EM Algorithm}
\label{sec:em-results}
% ===================================================================

We now evaluate the implemented EM algorithm
(\Cref{alg:weibull-em}) for Weibull parallel systems ($k = 1$),
comparing its performance against the direct MLE of
\Cref{sec:simulations}.  All experiments use $n = 300$ system
lifetime observations and $R = 200$ Monte Carlo replications
(unless noted otherwise), with the EM initialized by multiple
random starts ($n_{\text{starts}} = 3$) and the direct MLE using
L-BFGS-B with Nelder--Mead fallback ($n_{\text{starts}} = 5$).

\subsection{Shape effect on estimation difficulty}
\label{sec:em-shape-effect}

We first examine how the Weibull shape parameter affects estimation
difficulty for a two-component parallel system with scales
$(\beta_1, \beta_2) = (2, 3)$.  Both components share a common
shape $\alpha$ varied across $\{0.5, 1.0, 1.5, 2.0, 3.0\}$,
spanning DFR through strongly IFR regimes.

\begin{table}[t]
\centering
\caption{%
  Effect of shape parameter on EM and direct MLE for a
  two-component parallel system ($\bbeta = (2, 3)$,
  $n = 300$, $R = 200$).
  Shape RMSE and scale RMSE are averages over both components.
  The EM dominates in shape recovery, particularly at
  $\alpha = 1.0$ where the MLE's shape estimates explode
  for the fast component.
}
\label{tab:shape-effect}
\small
\begin{tabular}{c rr rr rr}
\toprule
& \multicolumn{3}{c}{EM} & \multicolumn{3}{c}{Direct MLE} \\
\cmidrule(lr){2-4} \cmidrule(lr){5-7}
$\alpha$ & Conv. & Shape & Scale
         & Conv. & Shape & Scale \\
         &  (\%) & RMSE & RMSE
         &  (\%) & RMSE & RMSE \\
\midrule
0.5 &  98 & 0.208 & 2.055 & 100 & 0.641 & 2.205 \\
1.0 & 100 & 0.315 & 0.819 & 100 & 7.487 & 0.912 \\
1.5 & 100 & 0.498 & 0.533 & 100 & 2.050 & 0.554 \\
2.0 & 100 & 0.738 & 0.400 & 100 & 0.914 & 0.403 \\
3.0 & 100 & 0.969 & 0.241 & 100 & 1.263 & 0.256 \\
\bottomrule
\end{tabular}
\end{table}

\Cref{tab:shape-effect} reveals two systematic patterns.
First, shape RMSE \emph{increases} with $\alpha$ (from 0.21 at
$\alpha = 0.5$ to 0.97 at $\alpha = 3$) while scale RMSE
\emph{decreases} (from 2.06 to 0.24).  Higher shapes concentrate
the Weibull distribution, improving scale estimation but creating
a steeper likelihood surface in which the shape is harder to pin
down.  Second, the EM dominates the direct MLE in shape recovery:
at $\alpha = 1.0$, the EM achieves shape RMSE $= 0.32$ while the
MLE explodes to 7.49 due to boundary violations in the fast
component.  Both methods achieve near-identical scale RMSE, as the
scale is better-identified than the shape under Scheme~0.  The EM's
advantage is clearest for the fast (deeply censored) component,
where the direct MLE's shape estimates are driven to extreme values
by the flat likelihood ridge along the shape--scale confounding
direction.

\subsection{Mixed shapes: DFR + IFR}
\label{sec:em-mixed}

In practice, systems contain components with heterogeneous
failure mechanisms.  We test three scenarios with mixed shapes:
\begin{itemize}
\item \textbf{Mild}: $\balpha = (0.8, 1.5)$,
  $\bbeta = (2.0, 3.0)$;
\item \textbf{Strong}: $\balpha = (0.5, 3.0)$,
  $\bbeta = (2.0, 3.0)$;
\item \textbf{Three-component}: $\balpha = (0.7, 1.5, 2.5)$,
  $\bbeta = (1.5, 2.0, 3.0)$.
\end{itemize}

The strong scenario is deliberately challenging: one DFR
component (shape $0.5$) with unbounded hazard at $t = 0$ and
one strongly IFR component (shape $3.0$) that fails in a narrow
window.  This creates a bimodal posterior over which component
failed last and stresses both the E-step classification
probabilities and the M-step profile optimization.

Both EM and MLE achieve 100\% convergence on the two-component
scenarios, with mean shape RMSE of 0.34 (mild) and 0.39
(strong).  The three-component case is harder: the fast
component ($\alpha = 0.7$, $\beta = 1.5$) has shape RMSE 3.60
under EM---the same information asymmetry seen earlier---while
the two slower components achieve RMSE below 0.85.  The EM
outperforms direct MLE in the three-component case: mean shape
RMSE 1.62 vs.\ 5.32, with the MLE's fast-component shape RMSE
reaching 14.7 due to boundary violations.

\subsection{EM vs direct MLE: convergence comparison}
\label{sec:em-convergence}

The results of \Cref{tab:shape-effect} and the mixed-shape
experiments reveal a consistent pattern: the EM achieves
comparable or superior estimation quality to the direct MLE,
with its advantage concentrated in \emph{shape} recovery.
At $\alpha = 1.0$, the EM's shape RMSE is 24 times smaller
(0.32 vs.\ 7.49); even at $\alpha = 2.0$, the gap persists
(0.74 vs.\ 0.91).  For three-component systems with mixed
shapes, the EM achieves mean shape RMSE 1.62 compared to
the MLE's 5.32---a 3.3-fold improvement.

Both methods achieve $\geq 98\%$ convergence for
two-component systems.  The EM's advantage is
\emph{structural}: its monotone likelihood guarantee prevents
the boundary violations and parameter blow-up that plague
gradient-based optimization on the Weibull likelihood surface,
particularly when shapes are near~1 (the exponential boundary)
or when components have very different scales.
Both methods achieve similar scale RMSE when they converge,
confirming they find the same optima---the EM simply reaches
them more reliably.

\subsection{Sample size requirements}
\label{sec:em-sample-size}

How much data does Weibull parallel system estimation require?
We vary $n \in \{50, 100, 200, 500\}$ for a two-component system
with $\balpha = (1.5, 2.0)$ and $\bbeta = (2.0, 3.0)$.

The results show a characteristic pattern: the slow component
(larger scale, more likely to be the system maximum) achieves
shape RMSE $\leq 0.56$ even at $n = 50$, while the fast component
(deeply left-censored under Scheme~0) has shape RMSE 5.76 at
$n = 50$, improving to 2.50 at $n = 200$ and dropping below~1
only at $n = 500$ (RMSE $= 0.60$).  This asymmetric sample
size requirement directly reflects the information asymmetry
of the parallel system: the last-to-fail component is observed
exactly, while earlier-failing components contribute only through
their CDF at the system failure time.  Both EM and MLE achieve
$\geq 96\%$ convergence across all sample sizes.


% ===================================================================
\section{Empirical Results: Scheme~1}
\label{sec:scheme1-results}
% ===================================================================

The EM results of \Cref{sec:em-results} confirm that Scheme~0
estimation is feasible for Weibull parallel systems but suffers
from the information asymmetry identified in \citet{Towell2026}:
fast components are poorly estimated.  We now evaluate whether
periodic inspection (Scheme~1) resolves this difficulty.

\subsection{Scheme~0 vs Scheme~1 comparison}
\label{sec:scheme1-comparison}

We compare Scheme~0 (EM, system-level only) with Scheme~1
(MLE using system failure time plus interval-censored component
data) for a two-component Weibull system with
$\balpha = (1.5, 2.0)$ and $\bbeta = (2.0, 3.0)$, using
$n = 300$ observations and $R = 200$ replications.  The
inspection interval $\delta$ is varied across
$\{0.1, 0.5, 1.0, 2.0\}$.

\begin{table}[t]
\centering
\caption{%
  Scheme~0 (EM) vs Scheme~1 (MLE with periodic inspection)
  for a two-component Weibull parallel system
  ($\balpha = (1.5, 2.0)$, $\bbeta = (2.0, 3.0)$,
  $n = 300$, $R = 200$).  Scheme~1 reduces shape RMSE
  by 12--21$\times$ for the fast component (Comp.~1)
  while improving the slow component (Comp.~2) by 3$\times$.
  Scheme~0 convergence: 94\%; all Scheme~1 configurations: 100\%.
}
\label{tab:scheme-comparison}
\small
\begin{tabular}{l rrrr}
\toprule
& \multicolumn{2}{c}{Shape RMSE} & \multicolumn{2}{c}{Scale RMSE} \\
\cmidrule(lr){2-3} \cmidrule(lr){4-5}
Method & Comp.\ 1 & Comp.\ 2 & Comp.\ 1 & Comp.\ 2 \\
\midrule
Scheme~0 (EM)              & 2.042 & 0.347 & 0.791 & 0.337 \\
Scheme~1 ($\delta = 0.1$)  & 0.097 & 0.107 & 0.099 & 0.103 \\
Scheme~1 ($\delta = 0.5$)  & 0.080 & 0.111 & 0.103 & 0.102 \\
Scheme~1 ($\delta = 1.0$)  & 0.095 & 0.125 & 0.120 & 0.116 \\
Scheme~1 ($\delta = 2.0$)  & 0.166 & 0.109 & 0.110 & 0.097 \\
\bottomrule
\end{tabular}
\end{table}

\Cref{tab:scheme-comparison} confirms the central prediction:
Scheme~1 provides massive improvement for the fast component
(Component~1, $\beta_1 = 2$), whose shape RMSE drops from 2.04
under Scheme~0 to 0.10 under Scheme~1 with $\delta = 0.1$---a
21-fold improvement.  The slow component (Component~2,
$\beta_2 = 3$) improves from 0.35 to 0.11, a more modest
3-fold gain.  This asymmetry is exactly what the censoring
framework predicts: the fast component is deeply left-censored
under Scheme~0 (its CDF is near~1 at typical system failure
times), so periodic inspection's interval constraints provide
maximal new information.

Even coarse inspection ($\delta = 2.0$) reduces shape RMSE by
a factor of~12 for the fast component.  Remarkably, performance
is relatively insensitive to the inspection interval: $\delta$
varying over a 20-fold range (0.1 to 2.0) produces shape RMSE
between 0.08 and 0.17 for the fast component.  This robustness
is encouraging for practical deployment, where fine-grained
inspection may be costly.

\subsection{Information-theoretic interpretation}

The Scheme~1 information gain has a clean interpretation in
terms of the censoring framework of \citet{Towell2026}.  Under
Scheme~0, the fast component contributes to the likelihood
primarily through $F_j(t)$---the probability of having failed by
time $t$---which is close to 1 for the fast component and
therefore nearly flat in $(\alpha_j, \beta_j)$.  Under Scheme~1,
the interval constraint $T_j \in (\underline{t}_{ij},
\bar{t}_{ij}]$ provides a direct localization of the component
failure time, giving likelihood contributions that are sharply
peaked in both shape and scale.

The data confirm this: the fast component's shape RMSE drops from
2.04 to $\sim$0.10 (a 95\% reduction), while its scale RMSE
drops from 0.79 to $\sim$0.10 (an 87\% reduction).  The larger
proportional gain for shape is because, under Scheme~0, the scale
is partially identifiable from the system CDF (which constrains
the marginal failure timing), but the shape---which governs how
the hazard changes over time---is nearly invisible from a single
system failure time.  Interval censoring at two time points
immediately constrains the hazard trend.


% ===================================================================
\section{Discussion}
\label{sec:discussion}
% ===================================================================

\subsection{Summary of findings}

The extension from exponential to Weibull component estimation in
$k$-out-of-$n$ systems is not a smooth generalization but a phase
transition in difficulty.  We have identified four distinct
mechanisms and two computational remedies:

\begin{enumerate}
\item \textbf{Parameter proliferation}: doubling the parameter
  count from $m$ to $2m$ thins the Fisher information per parameter
  and increases the dimensionality of the optimization landscape.

\item \textbf{Loss of analytical structure}: the
  inclusion-exclusion expansion that enables closed-form
  exponential likelihoods does not extend to Weibull components,
  requiring numerical evaluation of all likelihood quantities.

\item \textbf{Likelihood surface pathology}: shape--scale
  confounding creates ridges and local optima that degrade
  optimizer convergence from near-perfect (exponential) to as low
  as 47\% (Weibull mixed shapes at $k = 4$) under direct MLE.

\item \textbf{Shape-dependent optimal design}: the optimal system
  structure for estimation depends on the shape parameter---a
  quantity the experimenter typically does not know in advance,
  creating a circular design problem.

\item \textbf{EM resolves convergence}: the EM algorithm of
  \Cref{sec:em-algorithm}, with its monotone likelihood guarantee
  and decomposition into $m$ independent sub-problems, achieves
  near-perfect convergence across all shape configurations
  tested (\Cref{sec:em-results}), substantially outperforming
  direct MLE for challenging DFR and mixed-shape scenarios.

\item \textbf{Scheme~1 resolves information asymmetry}: periodic
  inspection provides interval-censored component data that reduces
  shape RMSE by an order of magnitude
  (\Cref{sec:scheme1-results}), with the largest gains for the
  deeply-censored fast components that are nearly unestimable under
  Scheme~0.
\end{enumerate}

\subsection{Practical implications}

\paragraph{When is exponential sufficient?}
If the primary goal is component estimation and the shape
parameters are believed to be close to 1, the exponential model of
\citet{Towell2026} is strongly preferable.  It provides
10--30 times better RMSE, guaranteed convergence, closed-form EM
steps, and proven identifiability.  The Weibull model should be
used only when there is substantive evidence of non-constant hazard
rates, and even then, Scheme~1 (periodic inspection) is strongly
recommended.

\paragraph{System design for Weibull estimation.}
Our results show that the ``safe'' choice is a parallel system
($k = 1$) or a 2-out-of-$n$ system ($k = 2$).  These structures
achieve competitive RMSE across all three shape scenarios tested,
with convergence rates above 73\%.  Higher $k$ values are risky:
$k = 4$ achieves the best RMSE under constant hazard but only 47\%
convergence under mixed shapes.  Series systems ($k = n$) should
be avoided entirely for Weibull estimation.

\paragraph{The case for periodic inspection.}
Scheme~0 provides insufficient information for reliable Weibull
parameter recovery, particularly for the shape parameters.  Our
Scheme~1 simulation results confirm that periodic inspection at
moderate intervals dramatically improves estimation: shape RMSE
drops by an order of magnitude, and the deeply-censored fast
components that are nearly unestimable under Scheme~0 become
well-recovered.  The information gain is asymmetric---Scheme~1
primarily benefits the components that are hardest to estimate
under Scheme~0, while already-well-estimated components see
modest improvement.  Even coarse inspection intervals
($\delta \approx$ median component lifetime) suffice.

\subsection{Limitations}

\paragraph{Replication count.}
The $k$-out-of-$n$ Monte Carlo study (\Cref{sec:simulations})
uses $R = 15$ replications per cell, which suffices for the
10--30-fold RMSE differences observed but is too small for precise
confidence intervals.  The targeted EM and Scheme~1 studies
(\Cref{sec:em-results,sec:scheme1-results}) use $R = 200$
replications, providing more reliable estimates of bias, RMSE,
and convergence rates for the parallel system case.

\paragraph{Identifiability.}
A formal proof of Weibull parallel system identifiability
(analogous to \citealt{BasuGhosh1980} for the exponential case)
remains open.  The numerical evidence is consistent with
identifiability for all configurations tested, but edge cases
(e.g., two components with the same shape but different scales)
may present difficulties.

\subsection{Future work}

\paragraph{Optimal $k$ characterization.}
The shape-dependent optimal $k$ merits theoretical investigation.
Is there a simple relationship between the shape parameter vector
$\balpha$ and the optimal $k$ for estimation?  Can the Fisher
information matrix (computed numerically) predict the optimal $k$
without exhaustive simulation?

\paragraph{Bayesian and regularized approaches.}
The convergence difficulties under Scheme~0 suggest that informative
priors on the shape parameters (e.g.,
$\alpha_j \sim \mathrm{Gamma}(a, b)$ with moderate concentration
around 1) could regularize the problem.  The Bayesian approach of
\citet{BheringPereiraPolpo2014} is a natural starting point.

\paragraph{Non-parallel structures.}
The EM algorithm of \Cref{sec:em-algorithm} and the Scheme~1
likelihood of \Cref{sec:scheme1} are stated for general
$k$-out-of-$n$ systems, but empirically validated only for the
parallel case ($k = 1$).  Extending the simulation study to
$k > 1$ with the EM algorithm and comparing convergence with
the direct MLE across the full $(k, \balpha)$ grid is a natural
next step.

\paragraph{Optimal inspection interval.}
The Scheme~1 results demonstrate large information gains even at
moderate inspection intervals, but the optimal choice of $\delta$
(balancing monitoring cost against estimation precision) has not
been characterized.  A decision-theoretic analysis that trades off
inspection cost against the Fisher information gain would provide
practical guidance for monitoring system design.


% ===================================================================
% Bibliography
% ===================================================================

\bibliographystyle{plainnat}
\bibliography{references}

\end{document}
